%==============================================================
% BAB II SOLUSI DAN INOVASI YANG DITAWARKAN
%==============================================================

\section{Solusi dan Inovasi yang Ditawarkan}

\subsection{Deskripsi Solusi dan Alur Klinis Berjenjang}

TropiCare adalah aplikasi web skrining awal penyakit tropis dan penyakit jangka panjang (PTM) yang mudah diakses tanpa instalasi dan tanpa kewajiban membuat akun. Seluruh logika medis berpedoman pada Keputusan Menteri Kesehatan RI No. HK.01.07/MENKES/1186/2022 tentang Panduan Praktik Klinis bagi Dokter di FKTP \citep{Kemenkes2022a} serta Permenkes No. 3 Tahun 2026 tentang Penanggulangan Penyakit \citep{Permenkes2026}. Sistem bekerja dalam waktu sekitar 3 menit melalui empat lapisan penapisan:

\begin{enumerate}
    \item \textbf{Gerbang Privasi, Data Diri, \& Deteksi Cepat 11 Tanda Bahaya (\textit{Red Flags}):} 
    Setelah memilih jalur skrining (infeksi akut atau risiko jangka panjang) dan mencentang persetujuan privasi, pengguna mengisi data diri dasar (usia, jenis kelamin, dan status kehamilan) lalu sistem langsung memeriksa 11 kondisi gawat darurat: (1) muntah hebat atau tangan-kaki dingin (syok dengue), (2) perdarahan sulit berhenti atau batuk darah banyak, (3) tidak buang air kecil $>6\text{--}8$ jam atau urin gelap (gagal ginjal akut), (4) sesak napas sangat berat, (5) penurunan kesadaran atau kejang, (6) sangat lemas dan mata cekung (dehidrasi berat), (7) luka gigitan hewan penular rabies, (8) nyeri dada kiri menjalar (serangan jantung), (9) tanda stroke mendadak (SeGeRa Ke RS), (10) krisis darah tinggi, serta (11) tanda koma diabetes \citep{Kemenkes2022a}. Satu saja tanda bahaya terpilih (atau muncul mata/kulit sangat kuning saat demam pada gejala penyerta), kuesioner langsung berhenti dan layar berubah menjadi Rute Merah tenang tanpa animasi yang menampilkan tombol telepon 119, nomor darurat 119/112 yang dapat disalin, serta panduan pertolongan pertama (\textit{DO} dan \textit{DON'T}) sambil menuju rumah sakit. Deteksi ini dieksekusi di peramban ponsel dalam $<500\text{ ms}$ walau peladen sedang mati.

    \item \textbf{Kuesioner Adaptif, Peta Ekologi (\textit{One Health}), \& Peta Nyeri:} 
    Jika tidak ada tanda bahaya, pertanyaan menyesuaikan jawaban sebelumnya ke dalam 5 klaster infeksi tropis (demam, napas, pencernaan, kulit, bengkak) untuk 20 penyakit tropis atau 3 klaster penyakit kronis (gula \& saraf, tekanan \& kepala, sendi) untuk 6 kondisi PTM \citep{PERKENI2024}. Sistem menggali riwayat lingkungan 1--4 bulan terakhir (daerah endemis malaria di NTT/Papua atau skistosomiasis di Sulawesi Tengah, genangan banjir leptospirosis, kontak TBC, lingkungan padat skabies) \citep{Permenkes2026}, dilengkapi peta nyeri tubuh interaktif dan skala nyeri 0--10 berbasis wajah.

    \item \textbf{Pemeriksaan Keamanan Obat \& Proksi Fisik Tanpa Alat Medis:} 
    Saat pengguna mengetik atau memilih obat yang sudah diminum, sistem memeriksa keamanan obat (4 aturan intervensi aktif ditambah 3 skema kewaspadaan komorbid dengue). Suspek demam berdarah yang meminum NSAID (Ibuprofen, Aspirin, Asam Mefenamat) langsung mendapat peringatan kritis merah untuk menghentikannya demi mencegah perdarahan lambung \citep{Alam2022,Mohiuddin2023,Tovar2024}; penggunaan obat penahan diare (Loperamide) pada diare berdarah memicu peringatan kuning \citep{Koo2007,Schiller2020}; sedangkan bengkak kaki akibat Amlodipine dan batuk kering akibat Captopril diberi catatan informasi biru agar tidak disalahartikan sebagai filariasis atau TBC \citep{Liang2022,Savage2020,Sica2003,Dicpinigaitis2014}. Bagi rumah tangga tanpa termometer atau timbangan \citep{Kemenkes2024}, tersedia pertanyaan pengganti (menggigil hebat, nadi cepat saat istirahat, dan perubahan kelonggaran pakaian) serta input tanda vital yang otomatis menaikkan rute ke Kuning bila suhu $\ge 39,5^\circ\text{C}$ atau tensi sistolik $>160\text{ mmHg}$ \citep{PERKENI2024}.

    \item \textbf{Orkestrasi 4 Rute Layanan \& Tindak Lanjut Faskes:} 
    Hasil skrining menyajikan satu perkiraan utama beserta skor kecocokan (atau \textit{``belum dapat ditentukan''} yang diarahkan ke Rute Kuning bila skor $<30\%$), tiga kemungkinan lain, alasan \textit{``Kenapa aku menebak ini?''}, serta satu dari empat rute layanan:
    \begin{itemize}
        \item \textbf{Rute Merah (IGD Rumah Sakit \& 119):} Kondisi darurat diarahkan langsung ke IGD RS terdekat atau Rabies Center.
        \item \textbf{Rute Kuning Khusus (Poli Program Puskesmas):} Dugaan penyakit program nasional (TBC, Kusta, Frambusia, Filariasis, Skistosomiasis, atau Campak) diarahkan ke Puskesmas penyelenggara program.
        \item \textbf{Rute Kuning (Puskesmas / FKTP):} Perlu pemeriksaan dokter dan laboratorium; pada jalur kronis dilengkapi instruksi wajib puasa 10--12 jam (hanya minum air putih) sebelum cek gula darah dan kolesterol \citep{PERKENI2024}.
        \item \textbf{Rute Hijau (Rawat Mandiri):} Gejala ringan non-komplikasi dipandu perawatan mandiri di rumah dan apotek terdekat.
    \end{itemize}
\end{enumerate}

\subsection{Arsitektur Sistem dan Alur Data Transparan}

Alur kerja sistem dari perangkat pasien hingga layanan faskes dirancang secara transparan dengan memisahkan pemrosesan di perangkat, di peladen TropiCare, dan di layanan AI eksternal (Gambar~\ref{fig:arsitektur}).

\begin{figure}[H]
\centering
\resizebox{0.96\textwidth}{!}{%
\begin{tikzpicture}[
    node distance=0.55cm and 0.45cm,
    box/.style={rectangle, draw=black!70, fill=black!4, rounded corners=2mm, text centered, font=\footnotesize, minimum height=0.85cm, line width=0.7pt},
    cloud/.style={rectangle, draw=black, fill=black!10, rounded corners=2mm, text centered, font=\footnotesize\bfseries, minimum height=0.85cm, line width=0.8pt},
    arrow/.style={-{Stealth[scale=0.8]}, draw=black!70, line width=0.7pt}
]
    % Top pipeline
    \node[box, text width=4.3cm] (input) {\textbf{Di Perangkat Pasien (React)}\\Persetujuan Privasi, Lokasi $\rightarrow$ Provinsi, \& Cek 11 Tanda Bahaya};
    \node[cloud, text width=4.6cm, right=of input] (engine) {\textbf{Peladen TropiCare (FastAPI)}\\Endpoint Stateless, Aturan Rute, \& Mirror 13.930 Faskes};
    \node[box, text width=4.3cm, right=of engine] (output) {\textbf{Layar Hasil \& Tindak Lanjut}\\4 Rute, 5 Faskes Terdekat, Mobile JKN, \& QR Petugas};

    % Bottom pipeline
    \node[box, text width=4.3cm, below=0.7cm of input] (red) {\textbf{Layar MERAH + 119}\\Jalan langsung di browser walau peladen mati};
    \node[box, text width=4.6cm, below=0.7cm of engine] (ai) {\textbf{Layanan AI (Gemma 4 / MedGemma)}\\Skor kecocokan + RAG 203 kutipan \& kenali teks obat};
    \node[box, text width=4.3cm, below=0.7cm of output] (qr) {\textbf{PostgreSQL \& Pemindai QR}\\Statistik tanpa identitas langsung; rangkuman dibuka petugas};

    % Arrows
    \draw[arrow] (input) -- (engine);
    \draw[arrow] (engine) -- (output);
    \draw[arrow] (input) -- (red);
    \draw[arrow] (engine) -- (ai);
    \draw[arrow] (ai) -- (engine);
    \draw[arrow] (output) -- (qr);
\end{tikzpicture}%
}
\caption{Alur Sistem TropiCare dari Perangkat Pasien, Peladen FastAPI, hingga Layanan AI dan Petugas Faskes}
\label{fig:arsitektur}
\end{figure}

Keunggulan rekayasa produk TropiCare bertumpu pada lima pilar teknis:
\begin{enumerate}
    \item \textbf{Evaluasi Darurat di Peramban (\textit{Client-Side Safety Gate}):} Pemeriksaan 11 tanda bahaya darurat dan pembuatan kode QR berjalan langsung di peramban ponsel. Jika tanda bahaya terpilih, layar darurat Merah dan tombol 119 aktif seketika meskipun koneksi internet atau peladen sedang mati.
    \item \textbf{Privasi Terjaga (\textit{Zero-NIK} \& Lokasi di Perangkat):} Pengguna dapat menyelesaikan skrining sebagai tamu tanpa akun dan tanpa pernah dimintai NIK atau nama. Koordinat GPS hanya diolah di memori peramban bersama poligon batas wilayah untuk mengenali provinsi dan mengurutkan jarak 5 faskes terdekat; hanya nama provinsi yang dikirim ke peladen \citep{UUPDP2022}.
    \item \textbf{Mirror Lokal 13.930 Faskes \& Jembatan Mobile JKN:} Peladen memuat mirror lokal 13.930 faskes resmi SATUSEHAT (10.374 Puskesmas dan 3.556 Rumah Sakit). Daftar 5 faskes terdekat dilengkapi tautan Google Maps, tombol kirim rangkuman via WhatsApp (atau telepon), serta tombol pembuka aplikasi Mobile JKN untuk pendaftaran janji dokter.
    \item \textbf{Penalaran Gemma 4 / MedGemma + RAG yang Terikat Pedoman Dokter:} Jawaban dianalisis menggunakan model keluarga Gemma (\textbf{Gemma 4 / MedGemma}) atau penyedia kompatibel (OpenRouter, vLLM, maupun Gemini) yang diikat pada 203 kutipan literatur medis terverifikasi (KMK 1186/2022 dan pedoman penyakit tropis Kemenkes) melalui indeks leksikal BM25 pra-komputasi. AI juga membantu mengenali nama obat ketikan pasien, sedangkan keputusan rute darurat dan ambang vital tetap dikunci oleh aturan deterministik.
    \item \textbf{QR Rangkuman untuk Petugas Faskes:} Jawaban pasien dikodekan secara ringkas ke dalam kode QR di layar ponsel tanpa menyimpan identitas di peladen. Saat dipindai petugas Puskesmas/klinik, rangkuman kondisi, riwayat obat, dan saran pemeriksaan laboratorium langsung terbuka di peramban petugas tanpa perlu wawancara ulang dari nol.
\end{enumerate}

\subsection{Perbandingan dengan Alternatif yang Ada}

TropiCare dirancang untuk \textbf{melengkapi, bukan menggantikan} layanan resmi seperti Mobile JKN: setelah tahu harus ke mana dari hasil skrining TropiCare, pasien langsung menekan tombol Mobile JKN untuk mendaftar antrean dokter. Perbandingan fitur disajikan pada Tabel~\ref{tab:keunggulan}.

\begin{table}[H]
\centering
\footnotesize
\caption{Perbandingan TropiCare dengan Alternatif Aplikasi Kesehatan di Indonesia}\label{tab:keunggulan}
\vspace{3pt}
\begin{tabularx}{\textwidth}{p{3.6cm} X X X}
\toprule
\textbf{Aspek Pembanding} & \textbf{TropiCare} & \textbf{Mobile JKN (BPJS)} & \textbf{Ada Health} \\
\midrule
\textbf{Skrining Gejala Infeksi Tropis} & \textbf{Ya (20 penyakit tropis Indonesia)} & Tidak tersedia & Umum, tidak khusus penyakit tropis Indonesia \\
\midrule
\textbf{Skrining Risiko Kronis (PTM)} & \textbf{Ya (6 kondisi PTM + proksi fisik)} & Ya (47 pertanyaan tetap, dibatasi $1\times$/tahun) & Umum \\
\midrule
\textbf{Kuesioner Adaptif} & \textbf{Ya (menyesuaikan jawaban sebelumnya)} & Tidak (daftar pertanyaan tetap) & Ya \\
\midrule
\textbf{Tanda Bahaya $\rightarrow$ Darurat Langsung} & \textbf{Ya (jalan di browser walau server mati)} & Tidak tersedia & Saran umum menghubungi darurat \\
\midrule
\textbf{Arahan Jenis Layanan} & \textbf{4 Rute (IGD, Poli Program, Puskesmas, Mandiri)} & Tingkat risiko \& saran gaya hidup & Saran langkah berikutnya secara umum \\
\midrule
\textbf{Rangkuman Petugas (QR) \& Bahasa} & \textbf{QR rangkuman instan; Bahasa Indonesia} & Tanpa QR pra-triase; Bahasa Indonesia & Laporan asesmen; belum tersedia Bahasa Indonesia konsumen \\
\bottomrule
\end{tabularx}
\end{table}
